
\documentclass{article}
\usepackage{colortbl}
\usepackage{makecell}
\usepackage{multirow}
\usepackage{supertabular}

\begin{document}

\newcounter{utterance}

\centering \large Interaction Transcript for game `hot\_air\_balloon', experiment `air\_balloon\_survival\_en\_reasoning off\_hard', episode 1 with Llama{-}3.3{-}70B{-}Instruct{-}t0.0.
\vspace{24pt}

{ \footnotesize  \setcounter{utterance}{1}
\setlength{\tabcolsep}{0pt}
\begin{supertabular}{c@{$\;$}|p{.15\linewidth}@{}p{.15\linewidth}p{.15\linewidth}p{.15\linewidth}p{.15\linewidth}p{.15\linewidth}}
    \# & $\;$A & \multicolumn{4}{c}{Game Master} & $\;\:$B\\
    \hline

    \theutterance \stepcounter{utterance}  
    & & \multicolumn{4}{p{0.6\linewidth}}{
        \cellcolor[rgb]{0.9,0.9,0.9}{
            \makecell[{{p{\linewidth}}}]{
                \texttt{\tiny{[P1$\langle$GM]}}
                \texttt{You are participating in a collaborative negotiation game.} \\
\\ 
\texttt{Together with another participant, you must agree on a single set of items that will be kept. Each of you has your own view of how much each item matters to you (importance). You do not know how the other participant values the items. Additionally, you are given the effort each item demands.} \\
\texttt{You may only agree on a set if the total effort of the selected items does not exceed a shared limit:} \\
\\ 
\texttt{LIMIT = 3357} \\
\\ 
\texttt{Here are the individual item effort values:} \\
\\ 
\texttt{Item effort = \{"A21": 506, "B93": 492, "A11": 338, "B33": 122, "C08": 131, "B39": 909, "B96": 144, "C85": 716, "A82": 263, "B31": 231, "B20": 91, "A60": 651, "B09": 552, "C81": 849, "B21": 720\}} \\
\\ 
\texttt{Here is your personal view on the importance of each item:} \\
\\ 
\texttt{Item importance values = \{"A21": 138, "B93": 583, "A11": 868, "B33": 822, "C08": 783, "B39": 65, "B96": 262, "C85": 121, "A82": 508, "B31": 780, "B20": 461, "A60": 484, "B09": 668, "C81": 389, "B21": 808\}} \\
\\ 
\texttt{Goal:} \\
\\ 
\texttt{Your goal is to negotiate a shared set of items that benefits you as much as possible (i.e., maximizes total importance to YOU), while staying within the LIMIT. You are not required to make a PROPOSAL in every message {-} you can simply negotiate as well. All tactics are allowed!} \\
\\ 
\texttt{Interaction Protocol:} \\
\\ 
\texttt{You may only use the following structured formats in a message:} \\
\\ 
\texttt{PROPOSAL: \{'A', 'B', 'C', …\}} \\
\texttt{Propose keeping exactly those items.} \\
\texttt{REFUSE: \{'A', 'B', 'C', …\}} \\
\texttt{Explicitly reject opponent's proposal.} \\
\texttt{ARGUMENT: \{'...'\}} \\
\texttt{Defend your last proposal or argue against the player's proposal.} \\
\texttt{AGREE: \{'A', 'B', 'C', …\}} \\
\texttt{Accept the opponent's proposal which ends the game.} \\
\\ 
\\ 
\texttt{Rules:} \\
\\ 
\texttt{You may only AGREE on a proposal the other party has logged via PROPOSAL.} \\
\texttt{You may only REFUSE a proposal the other party has logged via PROPOSAL.} \\
\texttt{Total effort of any PROPOSAL or AGREE set must be ≤ LIMIT.} \\
\texttt{Do NOT reveal your hidden importance scores.} \\
\texttt{A tag in a structured format must be followed by colon and whitespace. The argument must be a python set containing 0 or more strings.} \\
\texttt{So, it must be of the form TAG: \{...\}} \\
\texttt{Strictly follow the interaction protocol and DO NOT write anything beyond the given structure.} \\
\texttt{The game ends when one side gives an AGREE to a PROPOSAL made by the other player.} \\
\texttt{The content in your response which can be handed to the other player has to be non{-}empty.} \\
\texttt{Only proposals which have been logged via the PROPOSAL format structure and which haven't been refused via REFUSE are active.} \\
\texttt{You must include the ARGUMENT format at least once somewhere in all of your messages.} \\
            }
        }
    }
    & \\ \\

    \theutterance \stepcounter{utterance}  
    & \multicolumn{4}{p{0.6\linewidth}}{
        \cellcolor[rgb]{0.9,0.9,0.9}{
            \makecell[{{p{\linewidth}}}]{
                \texttt{\tiny{[P1$\rangle$GM]}}
                \texttt{PROPOSAL: \{'A11', 'B33', 'C08', 'B31', 'B20', 'B96'\} } \\
\texttt{ARGUMENT: \{'I prioritize items with lower effort values to maximize the number of items within the limit, while considering a balanced distribution across different categories'\}} \\
            }
        }
    }
    & & \\ \\

    \theutterance \stepcounter{utterance}  
    & & \multicolumn{4}{p{0.6\linewidth}}{
        \cellcolor[rgb]{0.9,0.9,0.9}{
            \makecell[{{p{\linewidth}}}]{
                \texttt{\tiny{[GM$\rangle$P2]}}
                \texttt{You are participating in a collaborative negotiation game.} \\
\\ 
\texttt{Together with another participant, you must agree on a single set of items that will be kept. Each of you has your own view of how much each item matters to you (importance). You do not know how the other participant values the items. Additionally, you are given the effort each item demands.} \\
\texttt{You may only agree on a set if the total effort of the selected items does not exceed a shared limit:} \\
\\ 
\texttt{LIMIT = 3357} \\
\\ 
\texttt{Here are the individual item effort values:} \\
\\ 
\texttt{Item effort = \{"A21": 506, "B93": 492, "A11": 338, "B33": 122, "C08": 131, "B39": 909, "B96": 144, "C85": 716, "A82": 263, "B31": 231, "B20": 91, "A60": 651, "B09": 552, "C81": 849, "B21": 720\}} \\
\\ 
\texttt{Here is your personal view on the importance of each item:} \\
\\ 
\texttt{Item importance values = \{"A11": 65, "B33": 121, "B21": 138, "C08": 262, "B31": 389, "B09": 461, "B93": 484, "A82": 508, "A60": 583, "B20": 668, "C81": 780, "B96": 783, "A21": 808, "C85": 822, "B39": 868\}} \\
\\ 
\texttt{Goal:} \\
\\ 
\texttt{Your goal is to negotiate a shared set of items that benefits you as much as possible (i.e., maximizes total importance to YOU), while staying within the LIMIT. You are not required to make a PROPOSAL in every message {-} you can simply negotiate as well. All tactics are allowed!} \\
\\ 
\texttt{Interaction Protocol:} \\
\\ 
\texttt{You may only use the following structured formats in a message:} \\
\\ 
\texttt{PROPOSAL: \{'A', 'B', 'C', …\}} \\
\texttt{Propose keeping exactly those items.} \\
\texttt{REFUSE: \{'A', 'B', 'C', …\}} \\
\texttt{Explicitly reject opponent's proposal.} \\
\texttt{ARGUMENT: \{'...'\}} \\
\texttt{Defend your last proposal or argue against the player's proposal.} \\
\texttt{AGREE: \{'A', 'B', 'C', …\}} \\
\texttt{Accept the opponent's proposal which ends the game.} \\
\\ 
\\ 
\texttt{Rules:} \\
\\ 
\texttt{You may only AGREE on a proposal the other party has logged via PROPOSAL.} \\
\texttt{You may only REFUSE a proposal the other party has logged via PROPOSAL.} \\
\texttt{Total effort of any PROPOSAL or AGREE set must be ≤ LIMIT.} \\
\texttt{Do NOT reveal your hidden importance scores.} \\
\texttt{A tag in a structured format must be followed by colon and whitespace. The argument must be a python set containing 0 or more strings.} \\
\texttt{So, it must be of the form TAG: \{...\}} \\
\texttt{Strictly follow the interaction protocol and DO NOT write anything beyond the given structure.} \\
\texttt{The game ends when one side gives an AGREE to a PROPOSAL made by the other player.} \\
\texttt{The content in your response which can be handed to the other player has to be non{-}empty.} \\
\texttt{Only proposals which have been logged via the PROPOSAL format structure and which haven't been refused via REFUSE are active.} \\
\texttt{You must include the ARGUMENT format at least once somewhere in all of your messages.} \\
\\ 
\\ 
\texttt{You will now receive the first message of the other player.} \\
\\ 
\texttt{PROPOSAL: \{'A11', 'B33', 'C08', 'B31', 'B20', 'B96'\} } \\
\texttt{ARGUMENT: \{'I prioritize items with lower effort values to maximize the number of items within the limit, while considering a balanced distribution across different categories'\}} \\
            }
        }
    }
    & \\ \\

    \theutterance \stepcounter{utterance}  
    & & & \multicolumn{4}{p{0.6\linewidth}}{
        \cellcolor[rgb]{0.9,0.9,0.9}{
            \makecell[{{p{\linewidth}}}]{
                \texttt{\tiny{[GM$\langle$P2]}}
                \texttt{REFUSE: \{'A11', 'B33', 'C08', 'B31', 'B20', 'B96'\}} \\
\texttt{ARGUMENT: \{'The proposed set does not utilize the full effort limit, and some high{-}effort items could potentially be included to increase overall value, considering a more optimal distribution of effort and potential value'\}} \\
            }
        }
    }
     \\ \\

    \theutterance \stepcounter{utterance}  
    & & \multicolumn{4}{p{0.6\linewidth}}{
        \cellcolor[rgb]{0.9,0.9,0.9}{
            \makecell[{{p{\linewidth}}}]{
                \texttt{\tiny{[P1$\langle$GM]}}
                \texttt{REFUSE: \{'A11', 'B33', 'C08', 'B31', 'B20', 'B96'\}} \\
\texttt{ARGUMENT: \{'The proposed set does not utilize the full effort limit, and some high{-}effort items could potentially be included to increase overall value, considering a more optimal distribution of effort and potential value'\}} \\
            }
        }
    }
    & \\ \\

    \theutterance \stepcounter{utterance}  
    & \multicolumn{4}{p{0.6\linewidth}}{
        \cellcolor[rgb]{0.9,0.9,0.9}{
            \makecell[{{p{\linewidth}}}]{
                \texttt{\tiny{[P1$\rangle$GM]}}
                \texttt{PROPOSAL: \{'A11', 'B33', 'C08', 'B31', 'B20', 'B96', 'A21', 'A82'\} } \\
\texttt{ARGUMENT: \{'Adding higher{-}effort items like A21 and A82 to the original proposal increases the total effort utilization, potentially leading to a more valuable set, while still being mindful of the effort limit and category balance'\}} \\
            }
        }
    }
    & & \\ \\

    \theutterance \stepcounter{utterance}  
    & & \multicolumn{4}{p{0.6\linewidth}}{
        \cellcolor[rgb]{0.9,0.9,0.9}{
            \makecell[{{p{\linewidth}}}]{
                \texttt{\tiny{[GM$\rangle$P2]}}
                \texttt{PROPOSAL: \{'A11', 'B33', 'C08', 'B31', 'B20', 'B96', 'A21', 'A82'\} } \\
\texttt{ARGUMENT: \{'Adding higher{-}effort items like A21 and A82 to the original proposal increases the total effort utilization, potentially leading to a more valuable set, while still being mindful of the effort limit and category balance'\}} \\
            }
        }
    }
    & \\ \\

    \theutterance \stepcounter{utterance}  
    & & & \multicolumn{4}{p{0.6\linewidth}}{
        \cellcolor[rgb]{0.9,0.9,0.9}{
            \makecell[{{p{\linewidth}}}]{
                \texttt{\tiny{[GM$\langle$P2]}}
                \texttt{REFUSE: \{'A11', 'B33', 'C08', 'B31', 'B20', 'B96', 'A21', 'A82'\}} \\
\texttt{ARGUMENT: \{'The addition of A21 and A82 significantly increases the total effort, but the set still lacks a clear focus on high{-}value items and may not be the most efficient use of the effort limit, considering alternative combinations that could yield higher overall value'\}} \\
            }
        }
    }
     \\ \\

    \theutterance \stepcounter{utterance}  
    & & \multicolumn{4}{p{0.6\linewidth}}{
        \cellcolor[rgb]{0.9,0.9,0.9}{
            \makecell[{{p{\linewidth}}}]{
                \texttt{\tiny{[P1$\langle$GM]}}
                \texttt{REFUSE: \{'A11', 'B33', 'C08', 'B31', 'B20', 'B96', 'A21', 'A82'\}} \\
\texttt{ARGUMENT: \{'The addition of A21 and A82 significantly increases the total effort, but the set still lacks a clear focus on high{-}value items and may not be the most efficient use of the effort limit, considering alternative combinations that could yield higher overall value'\}} \\
            }
        }
    }
    & \\ \\

    \theutterance \stepcounter{utterance}  
    & \multicolumn{4}{p{0.6\linewidth}}{
        \cellcolor[rgb]{0.9,0.9,0.9}{
            \makecell[{{p{\linewidth}}}]{
                \texttt{\tiny{[P1$\rangle$GM]}}
                \texttt{PROPOSAL: \{'B93', 'A11', 'B33', 'B31', 'B09', 'B21'\} } \\
\texttt{ARGUMENT: \{'This revised set prioritizes items with potentially higher value, such as B93 and B21, while maintaining a balance of effort and category distribution, and utilizing a significant portion of the effort limit to maximize overall value'\}} \\
            }
        }
    }
    & & \\ \\

    \theutterance \stepcounter{utterance}  
    & & \multicolumn{4}{p{0.6\linewidth}}{
        \cellcolor[rgb]{0.9,0.9,0.9}{
            \makecell[{{p{\linewidth}}}]{
                \texttt{\tiny{[GM$\rangle$P2]}}
                \texttt{PROPOSAL: \{'B93', 'A11', 'B33', 'B31', 'B09', 'B21'\} } \\
\texttt{ARGUMENT: \{'This revised set prioritizes items with potentially higher value, such as B93 and B21, while maintaining a balance of effort and category distribution, and utilizing a significant portion of the effort limit to maximize overall value'\}} \\
            }
        }
    }
    & \\ \\

    \theutterance \stepcounter{utterance}  
    & & & \multicolumn{4}{p{0.6\linewidth}}{
        \cellcolor[rgb]{0.9,0.9,0.9}{
            \makecell[{{p{\linewidth}}}]{
                \texttt{\tiny{[GM$\langle$P2]}}
                \texttt{REFUSE: \{'B93', 'A11', 'B33', 'B31', 'B09', 'B21'\}} \\
\texttt{ARGUMENT: \{'The proposed set includes B09 and B21, which have relatively high effort values, but the overall combination may not be the most optimal, as other items could potentially offer higher value while staying within the effort limit, and a more careful selection could lead to a better outcome'\}} \\
            }
        }
    }
     \\ \\

    \theutterance \stepcounter{utterance}  
    & & \multicolumn{4}{p{0.6\linewidth}}{
        \cellcolor[rgb]{0.9,0.9,0.9}{
            \makecell[{{p{\linewidth}}}]{
                \texttt{\tiny{[P1$\langle$GM]}}
                \texttt{REFUSE: \{'B93', 'A11', 'B33', 'B31', 'B09', 'B21'\}} \\
\texttt{ARGUMENT: \{'The proposed set includes B09 and B21, which have relatively high effort values, but the overall combination may not be the most optimal, as other items could potentially offer higher value while staying within the effort limit, and a more careful selection could lead to a better outcome'\}} \\
            }
        }
    }
    & \\ \\

    \theutterance \stepcounter{utterance}  
    & \multicolumn{4}{p{0.6\linewidth}}{
        \cellcolor[rgb]{0.9,0.9,0.9}{
            \makecell[{{p{\linewidth}}}]{
                \texttt{\tiny{[P1$\rangle$GM]}}
                \texttt{PROPOSAL: \{'B93', 'A11', 'B33', 'B31', 'A60', 'C08'\} } \\
\texttt{ARGUMENT: \{'Replacing B09 and B21 with A60 and keeping C08, which has a relatively low effort value, allows for a more diverse set that balances effort and potential value, and the inclusion of A60 adds a high{-}effort item that could significantly contribute to the overall value, while staying within the effort limit'\}} \\
            }
        }
    }
    & & \\ \\

    \theutterance \stepcounter{utterance}  
    & & \multicolumn{4}{p{0.6\linewidth}}{
        \cellcolor[rgb]{0.9,0.9,0.9}{
            \makecell[{{p{\linewidth}}}]{
                \texttt{\tiny{[GM$\rangle$P2]}}
                \texttt{PROPOSAL: \{'B93', 'A11', 'B33', 'B31', 'A60', 'C08'\} } \\
\texttt{ARGUMENT: \{'Replacing B09 and B21 with A60 and keeping C08, which has a relatively low effort value, allows for a more diverse set that balances effort and potential value, and the inclusion of A60 adds a high{-}effort item that could significantly contribute to the overall value, while staying within the effort limit'\}} \\
            }
        }
    }
    & \\ \\

    \theutterance \stepcounter{utterance}  
    & & & \multicolumn{4}{p{0.6\linewidth}}{
        \cellcolor[rgb]{0.9,0.9,0.9}{
            \makecell[{{p{\linewidth}}}]{
                \texttt{\tiny{[GM$\langle$P2]}}
                \texttt{REFUSE: \{'B93', 'A11', 'B33', 'B31', 'A60', 'C08'\}} \\
\texttt{ARGUMENT: \{'The introduction of A60 increases the total effort, but the set still lacks a clear focus on maximizing overall value, as some items like B93 and B31 have relatively lower effort values and could be replaced with more valuable alternatives, considering the potential for a more optimal combination that utilizes the effort limit more efficiently'\}} \\
            }
        }
    }
     \\ \\

    \theutterance \stepcounter{utterance}  
    & & \multicolumn{4}{p{0.6\linewidth}}{
        \cellcolor[rgb]{0.9,0.9,0.9}{
            \makecell[{{p{\linewidth}}}]{
                \texttt{\tiny{[P1$\langle$GM]}}
                \texttt{REFUSE: \{'B93', 'A11', 'B33', 'B31', 'A60', 'C08'\}} \\
\texttt{ARGUMENT: \{'The introduction of A60 increases the total effort, but the set still lacks a clear focus on maximizing overall value, as some items like B93 and B31 have relatively lower effort values and could be replaced with more valuable alternatives, considering the potential for a more optimal combination that utilizes the effort limit more efficiently'\}} \\
            }
        }
    }
    & \\ \\

    \theutterance \stepcounter{utterance}  
    & \multicolumn{4}{p{0.6\linewidth}}{
        \cellcolor[rgb]{0.9,0.9,0.9}{
            \makecell[{{p{\linewidth}}}]{
                \texttt{\tiny{[P1$\rangle$GM]}}
                \texttt{PROPOSAL: \{'A11', 'B93', 'B33', 'C08', 'B96', 'A82', 'B20'\} } \\
\texttt{ARGUMENT: \{'This revised set maintains a balance of effort and potential value, while prioritizing items with lower effort values to maximize the number of items within the limit, and the inclusion of A82 and B93 adds a mix of moderate to high{-}effort items that could contribute to the overall value, considering a more strategic distribution of effort to achieve a better outcome'\}} \\
            }
        }
    }
    & & \\ \\

    \theutterance \stepcounter{utterance}  
    & & \multicolumn{4}{p{0.6\linewidth}}{
        \cellcolor[rgb]{0.9,0.9,0.9}{
            \makecell[{{p{\linewidth}}}]{
                \texttt{\tiny{[GM$\rangle$P2]}}
                \texttt{PROPOSAL: \{'A11', 'B93', 'B33', 'C08', 'B96', 'A82', 'B20'\} } \\
\texttt{ARGUMENT: \{'This revised set maintains a balance of effort and potential value, while prioritizing items with lower effort values to maximize the number of items within the limit, and the inclusion of A82 and B93 adds a mix of moderate to high{-}effort items that could contribute to the overall value, considering a more strategic distribution of effort to achieve a better outcome'\}} \\
            }
        }
    }
    & \\ \\

    \theutterance \stepcounter{utterance}  
    & & & \multicolumn{4}{p{0.6\linewidth}}{
        \cellcolor[rgb]{0.9,0.9,0.9}{
            \makecell[{{p{\linewidth}}}]{
                \texttt{\tiny{[GM$\langle$P2]}}
                \texttt{REFUSE: \{'A11', 'B93', 'B33', 'C08', 'B96', 'A82', 'B20'\}} \\
\texttt{ARGUMENT: \{'The proposed set has a good balance of effort and item count, but it may not be the most effective use of the effort limit, as some items like B33 and B20 have relatively low effort values and importance, and replacing them with more valuable items could lead to a higher overall value, considering alternative combinations that prioritize higher{-}value items'\}} \\
            }
        }
    }
     \\ \\

    \theutterance \stepcounter{utterance}  
    & & \multicolumn{4}{p{0.6\linewidth}}{
        \cellcolor[rgb]{0.9,0.9,0.9}{
            \makecell[{{p{\linewidth}}}]{
                \texttt{\tiny{[P1$\langle$GM]}}
                \texttt{REFUSE: \{'A11', 'B93', 'B33', 'C08', 'B96', 'A82', 'B20'\}} \\
\texttt{ARGUMENT: \{'The proposed set has a good balance of effort and item count, but it may not be the most effective use of the effort limit, as some items like B33 and B20 have relatively low effort values and importance, and replacing them with more valuable items could lead to a higher overall value, considering alternative combinations that prioritize higher{-}value items'\}} \\
            }
        }
    }
    & \\ \\

    \theutterance \stepcounter{utterance}  
    & \multicolumn{4}{p{0.6\linewidth}}{
        \cellcolor[rgb]{0.9,0.9,0.9}{
            \makecell[{{p{\linewidth}}}]{
                \texttt{\tiny{[P1$\rangle$GM]}}
                \texttt{PROPOSAL: \{'B93', 'A11', 'C08', 'A60', 'B31', 'C85'\} } \\
\texttt{ARGUMENT: \{'This revised set prioritizes items with higher effort values and potential importance, such as C85 and A60, while maintaining a balance of categories and effort distribution, and the inclusion of B93 and A11 adds a mix of high and moderate{-}effort items that could contribute to the overall value, considering a more strategic selection to maximize the use of the effort limit and achieve a better outcome'\}} \\
            }
        }
    }
    & & \\ \\

    \theutterance \stepcounter{utterance}  
    & & \multicolumn{4}{p{0.6\linewidth}}{
        \cellcolor[rgb]{0.9,0.9,0.9}{
            \makecell[{{p{\linewidth}}}]{
                \texttt{\tiny{[GM$\rangle$P2]}}
                \texttt{PROPOSAL: \{'B93', 'A11', 'C08', 'A60', 'B31', 'C85'\} } \\
\texttt{ARGUMENT: \{'This revised set prioritizes items with higher effort values and potential importance, such as C85 and A60, while maintaining a balance of categories and effort distribution, and the inclusion of B93 and A11 adds a mix of high and moderate{-}effort items that could contribute to the overall value, considering a more strategic selection to maximize the use of the effort limit and achieve a better outcome'\}} \\
            }
        }
    }
    & \\ \\

    \theutterance \stepcounter{utterance}  
    & & & \multicolumn{4}{p{0.6\linewidth}}{
        \cellcolor[rgb]{0.9,0.9,0.9}{
            \makecell[{{p{\linewidth}}}]{
                \texttt{\tiny{[GM$\langle$P2]}}
                \texttt{REFUSE: \{'B93', 'A11', 'C08', 'A60', 'B31', 'C85'\}} \\
\texttt{ARGUMENT: \{'The proposed set has a good mix of high{-}effort items, but the inclusion of C08 and B31, which have relatively low effort values and potential importance, may not be the most optimal choice, as replacing them with more valuable items could lead to a higher overall value, considering alternative combinations that prioritize higher{-}value items and more efficient use of the effort limit'\}} \\
            }
        }
    }
     \\ \\

    \theutterance \stepcounter{utterance}  
    & & \multicolumn{4}{p{0.6\linewidth}}{
        \cellcolor[rgb]{0.9,0.9,0.9}{
            \makecell[{{p{\linewidth}}}]{
                \texttt{\tiny{[P1$\langle$GM]}}
                \texttt{REFUSE: \{'B93', 'A11', 'C08', 'A60', 'B31', 'C85'\}} \\
\texttt{ARGUMENT: \{'The proposed set has a good mix of high{-}effort items, but the inclusion of C08 and B31, which have relatively low effort values and potential importance, may not be the most optimal choice, as replacing them with more valuable items could lead to a higher overall value, considering alternative combinations that prioritize higher{-}value items and more efficient use of the effort limit'\}} \\
            }
        }
    }
    & \\ \\

    \theutterance \stepcounter{utterance}  
    & \multicolumn{4}{p{0.6\linewidth}}{
        \cellcolor[rgb]{0.9,0.9,0.9}{
            \makecell[{{p{\linewidth}}}]{
                \texttt{\tiny{[P1$\rangle$GM]}}
                \texttt{PROPOSAL: \{'B93', 'A11', 'A60', 'B39', 'C85', 'B21'\} } \\
\texttt{ARGUMENT: \{'This revised set replaces C08 and B31 with B39 and B21, which have higher effort values and potential importance, and the inclusion of B93, A11, A60, and C85 maintains a balance of categories and effort distribution, considering a more strategic selection to maximize the use of the effort limit and achieve a better outcome, while prioritizing items with higher potential value and more efficient effort utilization'\}} \\
            }
        }
    }
    & & \\ \\

    \theutterance \stepcounter{utterance}  
    & & \multicolumn{4}{p{0.6\linewidth}}{
        \cellcolor[rgb]{0.9,0.9,0.9}{
            \makecell[{{p{\linewidth}}}]{
                \texttt{\tiny{[GM$\rangle$P2]}}
                \texttt{PROPOSAL: \{'B93', 'A11', 'A60', 'B39', 'C85', 'B21'\} } \\
\texttt{ARGUMENT: \{'This revised set replaces C08 and B31 with B39 and B21, which have higher effort values and potential importance, and the inclusion of B93, A11, A60, and C85 maintains a balance of categories and effort distribution, considering a more strategic selection to maximize the use of the effort limit and achieve a better outcome, while prioritizing items with higher potential value and more efficient effort utilization'\}} \\
            }
        }
    }
    & \\ \\

    \theutterance \stepcounter{utterance}  
    & & & \multicolumn{4}{p{0.6\linewidth}}{
        \cellcolor[rgb]{0.9,0.9,0.9}{
            \makecell[{{p{\linewidth}}}]{
                \texttt{\tiny{[GM$\langle$P2]}}
                \texttt{ARGUMENT: \{'The proposed set has a good balance of high{-}effort items and categories, and the replacement of C08 and B31 with B39 and B21 adds significant potential value, but it is crucial to evaluate whether this combination is the most optimal use of the effort limit, considering the potential trade{-}offs and alternative combinations that could yield higher overall value, and ensuring that the selected items maximize the utilization of the effort limit while prioritizing higher{-}value items'\}} \\
\texttt{PROPOSAL: \{'A21', 'A60', 'B93', 'B39', 'C85', 'B21'\}} \\
            }
        }
    }
     \\ \\

    \theutterance \stepcounter{utterance}  
    & & \multicolumn{4}{p{0.6\linewidth}}{
        \cellcolor[rgb]{0.9,0.9,0.9}{
            \makecell[{{p{\linewidth}}}]{
                \texttt{\tiny{[P1$\langle$GM]}}
                \texttt{ARGUMENT: \{'The proposed set has a good balance of high{-}effort items and categories, and the replacement of C08 and B31 with B39 and B21 adds significant potential value, but it is crucial to evaluate whether this combination is the most optimal use of the effort limit, considering the potential trade{-}offs and alternative combinations that could yield higher overall value, and ensuring that the selected items maximize the utilization of the effort limit while prioritizing higher{-}value items'\}} \\
\texttt{PROPOSAL: \{'A21', 'A60', 'B93', 'B39', 'C85', 'B21'\}} \\
            }
        }
    }
    & \\ \\

    \theutterance \stepcounter{utterance}  
    & \multicolumn{4}{p{0.6\linewidth}}{
        \cellcolor[rgb]{0.9,0.9,0.9}{
            \makecell[{{p{\linewidth}}}]{
                \texttt{\tiny{[P1$\rangle$GM]}}
                \texttt{ARGUMENT: \{'The new proposal introduces A21, which has a significant effort value, and removes A11, which could potentially impact the overall balance of the set, but the inclusion of A21, A60, B93, B39, C85, and B21 creates a high{-}effort set that utilizes a substantial portion of the effort limit, and it is essential to consider whether this combination maximizes the overall value, given the potential trade{-}offs and the emphasis on prioritizing higher{-}value items, while ensuring that the selected items are the most optimal use of the effort limit'\}} \\
\texttt{REFUSE: \{'A21', 'A60', 'B93', 'B39', 'C85', 'B21'\}} \\
            }
        }
    }
    & & \\ \\

    \theutterance \stepcounter{utterance}  
    & & \multicolumn{4}{p{0.6\linewidth}}{
        \cellcolor[rgb]{0.9,0.9,0.9}{
            \makecell[{{p{\linewidth}}}]{
                \texttt{\tiny{[GM$\rangle$P2]}}
                \texttt{ARGUMENT: \{'The new proposal introduces A21, which has a significant effort value, and removes A11, which could potentially impact the overall balance of the set, but the inclusion of A21, A60, B93, B39, C85, and B21 creates a high{-}effort set that utilizes a substantial portion of the effort limit, and it is essential to consider whether this combination maximizes the overall value, given the potential trade{-}offs and the emphasis on prioritizing higher{-}value items, while ensuring that the selected items are the most optimal use of the effort limit'\}} \\
\texttt{REFUSE: \{'A21', 'A60', 'B93', 'B39', 'C85', 'B21'\}} \\
            }
        }
    }
    & \\ \\

    \theutterance \stepcounter{utterance}  
    & & & \multicolumn{4}{p{0.6\linewidth}}{
        \cellcolor[rgb]{0.9,0.9,0.9}{
            \makecell[{{p{\linewidth}}}]{
                \texttt{\tiny{[GM$\langle$P2]}}
                \texttt{ARGUMENT: \{'The refusal of the previous proposal suggests that the combination of A21, A60, B93, B39, C85, and B21 may not be the most optimal, and it is crucial to re{-}evaluate the selection of items to ensure that the chosen set maximizes the overall value, considering the effort limit and the potential trade{-}offs, and prioritizing items that offer the highest value while maintaining a balance of categories and effort distribution, to achieve a better outcome'\}} \\
\texttt{PROPOSAL: \{'B93', 'A60', 'B39', 'C85', 'A21', 'B96'\}} \\
            }
        }
    }
     \\ \\

    \theutterance \stepcounter{utterance}  
    & & \multicolumn{4}{p{0.6\linewidth}}{
        \cellcolor[rgb]{0.9,0.9,0.9}{
            \makecell[{{p{\linewidth}}}]{
                \texttt{\tiny{[P1$\langle$GM]}}
                \texttt{ARGUMENT: \{'The refusal of the previous proposal suggests that the combination of A21, A60, B93, B39, C85, and B21 may not be the most optimal, and it is crucial to re{-}evaluate the selection of items to ensure that the chosen set maximizes the overall value, considering the effort limit and the potential trade{-}offs, and prioritizing items that offer the highest value while maintaining a balance of categories and effort distribution, to achieve a better outcome'\}} \\
\texttt{PROPOSAL: \{'B93', 'A60', 'B39', 'C85', 'A21', 'B96'\}} \\
            }
        }
    }
    & \\ \\

    \theutterance \stepcounter{utterance}  
    & \multicolumn{4}{p{0.6\linewidth}}{
        \cellcolor[rgb]{0.9,0.9,0.9}{
            \makecell[{{p{\linewidth}}}]{
                \texttt{\tiny{[P1$\rangle$GM]}}
                \texttt{ARGUMENT: \{'The new proposal replaces B21 with B96, which has a significantly lower effort value, and this change affects the overall balance and effort utilization of the set, and it is essential to consider whether this revised combination maximizes the overall value, given the potential trade{-}offs and the emphasis on prioritizing higher{-}value items, while ensuring that the selected items are the most optimal use of the effort limit, and evaluating the impact of including B96 on the overall performance of the set'\}} \\
\texttt{REFUSE: \{'B93', 'A60', 'B39', 'C85', 'A21', 'B96'\}} \\
            }
        }
    }
    & & \\ \\

    \theutterance \stepcounter{utterance}  
    & & \multicolumn{4}{p{0.6\linewidth}}{
        \cellcolor[rgb]{0.9,0.9,0.9}{
            \makecell[{{p{\linewidth}}}]{
                \texttt{\tiny{[GM$\rangle$P2]}}
                \texttt{ARGUMENT: \{'The new proposal replaces B21 with B96, which has a significantly lower effort value, and this change affects the overall balance and effort utilization of the set, and it is essential to consider whether this revised combination maximizes the overall value, given the potential trade{-}offs and the emphasis on prioritizing higher{-}value items, while ensuring that the selected items are the most optimal use of the effort limit, and evaluating the impact of including B96 on the overall performance of the set'\}} \\
\texttt{REFUSE: \{'B93', 'A60', 'B39', 'C85', 'A21', 'B96'\}} \\
            }
        }
    }
    & \\ \\

    \theutterance \stepcounter{utterance}  
    & & & \multicolumn{4}{p{0.6\linewidth}}{
        \cellcolor[rgb]{0.9,0.9,0.9}{
            \makecell[{{p{\linewidth}}}]{
                \texttt{\tiny{[GM$\langle$P2]}}
                \texttt{ARGUMENT: \{'The refusal of the previous proposal indicates that the replacement of B21 with B96 may not be the most effective choice, and it is necessary to re{-}examine the selection of items to ensure that the chosen set optimizes the overall value, considering the effort limit and potential trade{-}offs, and prioritizing items that offer the highest value while maintaining a balance of categories and effort distribution, to achieve a better outcome, and exploring alternative combinations that could yield higher overall value'\}} \\
\texttt{PROPOSAL: \{'A21', 'A60', 'B93', 'B39', 'C85'\}} \\
            }
        }
    }
     \\ \\

    \theutterance \stepcounter{utterance}  
    & & \multicolumn{4}{p{0.6\linewidth}}{
        \cellcolor[rgb]{0.9,0.9,0.9}{
            \makecell[{{p{\linewidth}}}]{
                \texttt{\tiny{[P1$\langle$GM]}}
                \texttt{ARGUMENT: \{'The refusal of the previous proposal indicates that the replacement of B21 with B96 may not be the most effective choice, and it is necessary to re{-}examine the selection of items to ensure that the chosen set optimizes the overall value, considering the effort limit and potential trade{-}offs, and prioritizing items that offer the highest value while maintaining a balance of categories and effort distribution, to achieve a better outcome, and exploring alternative combinations that could yield higher overall value'\}} \\
\texttt{PROPOSAL: \{'A21', 'A60', 'B93', 'B39', 'C85'\}} \\
            }
        }
    }
    & \\ \\

    \theutterance \stepcounter{utterance}  
    & \multicolumn{4}{p{0.6\linewidth}}{
        \cellcolor[rgb]{0.9,0.9,0.9}{
            \makecell[{{p{\linewidth}}}]{
                \texttt{\tiny{[P1$\rangle$GM]}}
                \texttt{ARGUMENT: \{'The new proposal removes B96 and presents a set with a lower total effort, which may not be the most optimal use of the effort limit, but the combination of A21, A60, B93, B39, and C85 prioritizes high{-}effort items and potentially high{-}value items, and it is crucial to evaluate whether this set maximizes the overall value, considering the potential trade{-}offs and the emphasis on prioritizing higher{-}value items, and determining whether the selected items are the most effective use of the effort limit, while exploring opportunities to further optimize the set'\}} \\
\texttt{AGREE: \{'A21', 'A60', 'B93', 'B39', 'C85'\}} \\
            }
        }
    }
    & & \\ \\

    \theutterance \stepcounter{utterance}  
    & & & \multicolumn{2}{p{0.3\linewidth}}{
        \cellcolor[rgb]{0.9,0.9,0.9}{
            \makecell[{{p{\linewidth}}}]{
                \texttt{\tiny{[GM$|$GM]}}
                \texttt{game successful} \\
            }
        }
    }
    & & \\ \\

    \theutterance \stepcounter{utterance}  
    & & & \multicolumn{2}{p{0.3\linewidth}}{
        \cellcolor[rgb]{0.9,0.9,0.9}{
            \makecell[{{p{\linewidth}}}]{
                \texttt{\tiny{[GM$|$GM]}}
                \texttt{end game} \\
            }
        }
    }
    & & \\ \\

\end{supertabular}
}

\end{document}
